\section{Kuvani --- ``The Remount Steppe''}
\label{chap:kuvani}

\subsection*{Quick Reference}
\textbf{Tagline:} \textit{The sky hears every oath. The grass remembers every step.}

\begin{quote}
  \textbf{Wind-Singer (Elite):}  
  ``Kuva watches from the blue. He does not speak. He listens. A promise made under open sky is carved into the wind itself. The grass will bend around it. The horses will remember it. Break it, and the sky will turn its face away. You will ride forever in shadow.''
\end{quote}

\begin{quote}
  \textbf{Remount Master (Commoner):}  
  ``A horse's name is its lineage. Its lineage is its story. You do not ride a Kuvani horse. You borrow it. And when you return it, you will tell me a story that the horse has not already heard. If you lie, the horse will throw you next time — and Kuva will laugh.''
\end{quote}

\textbf{Genre / Mood:} Sky‑worshipping nomads, ancestral horse‑magic, wind‑carried oaths.

\textbf{Starting Location:} A felt tent circle at a remount station, where a corral of horses stands behind a rope of braided hair. A Remount Master sits on a saddle, a bell in his hand: ``You need fresh horses. I have twelve. Each horse has a name. Each name has a price. The price is not silver — it is a story I have not heard. Tell me one. If it makes me laugh, you ride for free. If it makes me weep, you owe me a foal. If Kuva sends a hawk to watch, you owe a prayer.''

\vspace{2mm}
\begin{tcolorbox}[colback=black!3,colframe=blue!60!black,title={Theme \& Atmosphere}]
The Kuvani ride the eastern steppe beneath the Eternal Blue Sky. Kuva, the sky father, sees all: every oath sworn in the open, every horse born, every blade of grass bent by the wind. The Kuvani have no temples — the sky is their shrine. Their shamans, called Wind-Singers, read the future in the flight of hawks and the whisper of grass. Their dead are not buried; they are given to the sky on funeral platforms, that Kuva may carry their souls to the ancestral herd that rides the northern lights.

Unlike their Ykrul neighbours, who worship geometry and oath‑strings, the Kuvani worship movement itself. A horse is not a mount — it is a relative. Its name is a lineage stretching back to the first mare that Kuva breathed life into. A Kuvani who cannot name his horse's ancestors is a Kuvani without a shadow. The steppe is not empty; it is full of ghosts: the ghost herd of the salt pan, the sky‑road markers that hum at dusk, the Iron Corral where blades whisper the names of the dead. A Kuvani does not fight for land — he fights for grazing rights, for the memory of grass, for the honour of his horse's line.

\vspace{2mm}
\textbf{What you notice first:} The smell of horse sweat and dry grass. The way every rider touches a small felt bundle at their belt — a piece of a dead horse's mane, a prayer to Kuva. The constant sound of bells on harnesses, each tone a different ancestor's name.

\textbf{The one rule that kills newcomers:} Never ride a Kuvani horse without asking its name. The horse will answer — by throwing you. If it throws you, you have insulted its lineage. The remount master will demand a story as compensation. If the story is not good enough, you walk — and Kuva will turn the wind against you.
\end{tcolorbox}

\subsection*{Regional Mood: The Weight of the Herd}

Power here is measured in remount chains, grazing compacts, and the favour of Kuva. A clan's status depends on the number of horses it can offer to caravans; a remount master's reputation on the speed and health of his string. The Kuvani have no kings — only Wind-Singers, who speak for the sky, and Remount Masters, who hold the earthly ledger of horse‑flesh. Every caravan that crosses the steppe must pay the remount toll: not in coin, but in stories. The Kuvani believe that a story shared is a prayer heard, and that a traveller who has no story has no soul — only an empty shape that the wind will soon erase.

\subsection*{Prominent NPCs of the Kuvani}
\label{sec:kuvani-npcs}

\begin{longtable}{|p{0.2\textwidth}|p{0.25\textwidth}|p{0.3\textwidth}|p{0.15\textwidth}|}
\hline
\textbf{Name} & \textbf{Role/Title} & \textbf{Motivation} & \textbf{Rumoured Quirk} \\
\hline
Alp Arslan & Remount master of the central route & Expand grazing rights into Ykrul territory & He has never been seen sleeping; his eyes are always on the horizon, watching for Kuva's hawks. \\
\hline
Borte & Wind-Singer of the White Mane Clan & Keep the songs of the clans from being forgotten & She has a horse with a white blaze; the blaze is a map of the steppe, and it shifts each season. \\
\hline
Temur & Caravan guide & Find the fastest routes through the shifting grass & He names his dogs after dead enemies. The dogs remember. They growl when the enemy's ghost is near. \\
\hline
Seljuk & Ykrul envoy (negotiator) & Negotiate shared pastures without war & He carries a hostage string with nine knots; each knot is a season of peace. He is running out of knots. \\
\hline
The Grass Reader & A hermit who lives in a felt hut & Read the future in the way the grass bends & He never speaks above a whisper. The grass speaks for him. His tongue is green. \\
\hline
\end{longtable}

\subsection*{Names of the Kuvani}
\label{sec:kuvani-names}

Inspired by Turkic and Mongolian traditions: short, hard consonants, vowel clusters, patronymic suffixes (-oglu, -ova). Surnames denote a clan, a horse-line, or a wind-name.

\begin{longtable}{|p{0.25\textwidth}|p{0.25\textwidth}|p{0.2\textwidth}|p{0.2\textwidth}|}
\hline
\textbf{First Names (Male)} & \textbf{First Names (Female)} & \textbf{Surnames} & \textbf{Epithets} \\
\hline
Alp & Borte & Arslanoglu & \textit{the Iron Rein} \\
\hline
Temur & Selin & Kuvani-bek & \textit{the White Mane} \\
\hline
Seljuk & Chagatai & Borjigin & \textit{Ninth Wind} \\
\hline
Batu & Nergui & Oirat & \textit{the Grass-Tongue} \\
\hline
Toghrul & Otgon & Kipchak & \textit{the Remount-Keeper} \\
\hline
\end{longtable}

\paragraph*{(Station/Circle/Kurgan)} Remount station with a felt tent and a corral; Wind-Singer's circle of saddles; a kurgan (burial mound) with a stone horse that is missing a leg.

\section*{Spades --- Places}
\label{sec:kuvani-places}
\begin{enumerate}
\setcounter{enumi}{1}
\item \textbf{Remount Station} --- A felt tent with a corral; horses for hire. The horses have no names — only brands.
  \hfill \textit{The brands are the names of lost children. The horses remember them. Kuva's hawks circle when a brand is new.} `[TRADE]`
\item \textbf{Wind-Singer's Circle} --- A ring of saddles; the singer sits in the centre. The songs are not sung — they are hummed.
  \hfill \textit{The hum is the wind. The wind carries the song to the next valley. The next valley answers with a hawk's cry.} `[RITUAL]`
\item \textbf{Kurgan of the Broken Oath} --- A burial mound with a stone horse missing a leg.
  \hfill \textit{The leg is a treaty. The treaty was broken. The horse is waiting for it to be returned. Its eye follows you.} `[DEATH]`
\item \textbf{Grass Maze} --- A section of the steppe where the grass grows taller than a rider on horseback.
  \hfill \textit{The paths shift with the wind. Only a Wind-Singer can find the way out. The grass whispers warnings.} `[MYSTERY]`
\item \textbf{Salt Pan of the Ghost Herd} --- A dry lake; the salt is crusted with hoofprints.
  \hfill \textit{The prints are from a herd that died a century ago. They appear every spring. The herd is still running. Kuva's hawks do not fly over it.} `[DEATH]`
\item \textbf{Felt Village (Mobile)} --- A temporary camp of round tents. The tents face the wind.
  \hfill \textit{The wind is the clan's direction. If the wind changes, the village moves. The door always faces the rising sun.} `[COMMUNITY]`
\item \textbf{Horse Cemetery} --- A field of skulls, each on a wooden stake. The stakes are carved with names.
  \hfill \textit{The names are the riders' names. The riders are buried elsewhere. The horses wait. Their skulls whicker at dawn.} `[DEATH]`
\item \textbf{Bridge of the Ninth Hardship} --- A log over a ravine. The log is replaced each year. The old log is burned.
  \hfill \textit{The ash is the clan's memory. The ash is scattered on the spring grass. Kuva drinks the smoke.} `[RITUAL]`
\item \textbf{Stone of the Hostage String} --- A standing stone with notches. Each notch is a hostage taken.
  \hfill \textit{The hostages are carved into the stone. The stone weeps when a hostage dies. The tears taste of salt and iron.} `[DEBT]`
\item[J] \textbf{The Iron Corral} --- A pen made of salvaged sword blades. The blades are from a war that never ended.
  \hfill \textit{The horses are afraid of them. The blades whisper the names of the dead. The whisper is the wind.} `[MAGIC]`
\item[Q] \textbf{Sky-Road Marker} --- A cairn painted blue. The blue is from a stone that fell from the sky.
  \hfill \textit{The stone is still warm. It marks the path of the summer stars. Touch it, and you will dream of falling.} `[TRAVEL]`
\item[K] \textbf{Khagan's Felt Throne} --- A tent within a tent; the seat is a saddle. The saddle has no stirrups.
  \hfill \textit{The Khagan does not need them. He rides the wind. The wind has no stirrups. The saddle was his grandfather's.} `[AUTHORITY]`
\item[A] \textbf{The Steppe's Edge} --- Where the grass ends and the desert begins. The line is a scar.
  \hfill \textit{The scar is a treaty. The treaty was signed in blood. The blood was Kuvani. Kuva turned his face away that day.} `[LEGEND]`
\end{enumerate}

\paragraph*{(Master/Singer/Guide)} Remount master with calloused hands and a horsehair whip; Wind-Singer with grey hair and a bone flute; caravan guide with a string of dog teeth.

\section*{Hearts --- People \& Factions}
\label{sec:kuvani-people}
\begin{enumerate}
\setcounter{enumi}{1}
\item \textbf{Remount Master} --- Calloused hands, squinting eyes, a whip of braided horsehair.
  \hfill \textit{He knows the price of a horse in salt, in silk, or in stories. The price changes with the wind. He has never been cheated.} `[TRADE]`
\item \textbf{Wind-Singer (Shaman)} --- Grey-haired, blind, but never lost. She hums the direction to the nearest water.
  \hfill \textit{The hum changes with the wind. If she stops humming, the wind stops too. Her left eye sees the living; her right sees the dead.} `[RITUAL]`
\item \textbf{Caravan Guide (Temur)} --- Young, fast, a scar across his lip. He names his dogs after dead enemies.
  \hfill \textit{The dogs remember. They growl when they smell the enemy's kin. The dogs also see ghosts.} `[TRAVEL]`
\item \textbf{Ykrul Envoy (Seljuk)} --- Braided hair, a hostage string on his wrist. The string has nine knots.
  \hfill \textit{Each knot is a season of peace. He is running out of knots. The last knot is tied with a hair from a Khagan's daughter.} `[POLITICS]`
\item \textbf{Grass Reader (Hermit)} --- A hermit in a felt hut, never seen without a handful of grass.
  \hfill \textit{He reads the future by how the grass bends. The future is always wind. He has never been wrong.} `[DIVINATION]`
\item \textbf{Horse Thief} --- Young, missing two fingers. The fingers were a bet. He lost the bet — but not the horses.
  \hfill \textit{He can steal a horse from a corral without waking the dogs. The dogs are his. He paid for them with a story.} `[STEALTH]`
\item \textbf{Felt Village Elder} --- A woman with a pipe made of bone. She settles disputes by listening to the wind.
  \hfill \textit{``The wind says you are lying,'' she says. The wind has never been wrong. The pipe belonged to her grandmother.} `[JUSTICE]`
\item \textbf{Khagan's Shadow} --- A silent figure in black felt. He never speaks; he writes commands in the dust.
  \hfill \textit{The dust blows away. The commands are remembered anyway. His face is never seen.} `[AUTHORITY]`
\item \textbf{Horse-Whisperer (Child)} --- A child who can calm any stallion. The child speaks in neighs.
  \hfill \textit{The horses answer in riddles. The child is learning to understand. One day, a horse will answer with a name.} `[MYSTERY]`
\item[J] \textbf{Ghost Herd's Rider} --- A translucent figure on a skeletal horse, seen at the salt pan at dawn.
  \hfill \textit{He asks: ``Where is my herd?'' If you answer, you become part of it. If you refuse, he follows you for nine days.} `[DEATH]`
\item[Q] \textbf{The Steppe's Voice (Oracle)} --- A rock with a crack. The wind speaks through it.
  \hfill \textit{The wind answers questions. The answers are always poems. The poems are always warnings. The rock is warm.} `[DIVINATION]`
\item[K] \textbf{The Last Khagan (Legend)} --- A king who united the clans, then rode into a storm and never returned.
  \hfill \textit{He is said to appear during the worst blizzards, offering shelter. The shelter is a trap. His horse has no shadow.} `[LEGEND]`
\item[A] \textbf{The First Remount (Myth)} --- A horse that never tired. Its bones are scattered across the steppe.
  \hfill \textit{Find all nine bones, and you can summon it. The summoning costs a year of your life. The horse will remember your name.} `[MYTH]`
\end{enumerate}

\paragraph*{(Sickness/Fire/Stampede)} A horse sickness spreads through the remounts — all mounts lose 1 Speed; grass fire — the smoke spells a name; a wild horse stampede — the horses are fleeing something worse.

\section*{Clubs --- Complications/Threats}
\label{sec:kuvani-complications}
\begin{enumerate}
\setcounter{enumi}{1}
\item \textbf{Horse Sickness (Curse)} --- A cough spreads through the remounts. The cough is a curse from a Ykrul bone-singer.
  \hfill \textit{The cure is a story that the bone-singer has not heard. The story must be true. Kuva's hawks will judge.} `[CURSE]`
\item \textbf{Grass Fire} --- The steppe burns; the smoke is black. The smoke spells a name.
  \hfill \textit{The name is the one who lit the fire. The wind will carry the name to the clans. They will hunt. The name also appears in your dreams.} `[DISASTER]`
\item \textbf{Stampede (Wild Herd)} --- A herd of wild horses crashes through camp. The horses are not wild.
  \hfill \textit{They are fleeing something worse. Something that hunts horses. The party is in its path. The ground shakes.} `[COMBAT]`
\item \textbf{Hostage String Cut} --- A Ykrul envoy's string is broken. The clan demands a new hostage.
  \hfill \textit{The party must choose a member to offer. The hostage will be held for nine seasons. Kuva will witness the oath.} `[DEBT]`
\item \textbf{Rasputitsa (Mud)} --- The spring thaw turns the steppe to mud. Wagons sink. Progress slows.
  \hfill \textit{The mud is warm. Something is buried underneath. It is waking up. The mud smells of old blood.} `[WEATHER]`
\item \textbf{Wind-Singer's Silence} --- The Wind-Singer stops humming. The wind stops too. The silence is a warning.
  \hfill \textit{The warning is: ``A storm is coming. Not of rain — of ghosts.'' The ghosts are the unburied.} `[DEATH]`
\item \textbf{Water Poisoned (Salt Pan)} --- A well is found with dead animals nearby. The poison is salt from the cursed pan.
  \hfill \textit{The salt pan wants a sacrifice. The sacrifice is a horse. The horse must be willing. The horse will neigh its acceptance.} `[CURSE]`
\item \textbf{Felt Village Gone} --- A camp is empty; the tents are still standing. The felt is warm.
  \hfill \textit{The people vanished at noon. The horses are still in the corral. The horses are trembling. Their eyes are white.} `[MYSTERY]`
\item \textbf{Khagan's Demand} --- The Khagan (or his shadow) asks for a tribute. The tribute is a story.
  \hfill \textit{The story must be about a broken promise. If it is not good enough, the Khagan takes a horse. The horse will never be seen again.} `[DEBT]`
\item[J] \textbf{Iron Corral Bleeding} --- The sword-blade fence begins to weep rust. The rust is blood.
  \hfill \textit{The blood is from the war that never ended. The war is about to resume. The blades are warm.} `[OMEN]`
\item[Q] \textbf{Sky-Road Marker Falls} --- The blue cairn collapses. The stone inside is a meteor. It whispers.
  \hfill \textit{The whisper is a prophecy. The prophecy is about the end of the steppe. The stone is cold.} `[OMEN]`
\item[K] \textbf{Ghost Herd Charges} --- Spectral horses trample the camp. They are looking for their riders.
  \hfill \textit{The riders are dead. The ghosts do not know. One rider's ghost is in your party. It does not know it is dead.} `[DEATH]`
\item[A] \textbf{Steppe's Edge Recedes} --- The grass grows into the desert. The desert is retreating.
  \hfill \textit{Something is driving the desert back. Something that wants the steppe. Kuva is silent.} `[APOCALYPSE]`
\end{enumerate}

\paragraph*{(Mark/Compact/Song)} Remount chain mark (a horse's tooth, exchange for a fresh mount); grazing compact (a sheepskin document, safe passage on clan land); sky-road song (a melody on a bone flute, navigate any storm).

\section*{Diamonds --- Rewards/Leverage}
\label{sec:kuvani-rewards}
\begin{enumerate}
\setcounter{enumi}{1}
\item \textbf{Remount Chain Mark} --- A horse's tooth on a leather thong. Exchange it for a fresh mount at any Kuvani station.
  \hfill \textit{The tooth is from a horse that died in battle. Its spirit runs with the new mount. The new mount will be braver.} `[TRAVEL]`
\item \textbf{Grazing Compact (Sheepskin)} --- A written agreement. Allows your caravan to pasture on clan land for one season.
  \hfill \textit{The compact is sealed with horsehair. The hair is from the Khagan's own mare. Break the seal, and the mare will know.} `[TRADE]`
\item \textbf{Sky-Road Song (Bone Flute)} --- A melody recorded on a flute. Play it to navigate any storm.
  \hfill \textit{The flute was carved from a rib. The rib was the Wind-Singer's. She gave it willingly. The flute is warm.} `[TRAVEL]`
\item \textbf{Hostage String (Borrowed)} --- A braided cord with knots. Present it to demand a truce from a Ykrul clan.
  \hfill \textit{The string is hair. The hair is from a Khagan's daughter. The Ykrul will recognise it. They will honour it. Once.} `[SOCIAL]`
\item \textbf{Felt Tent Pass} --- A felt circle with a hole. Safe lodging in any Kuvani camp.
  \hfill \textit{The hole is for smoke. The smoke is your story. The clan will listen. If they laugh, you stay. If they weep, you owe a foal.} `[SOCIAL]`
\item \textbf{Kurgan Key (Stone Tablet)} --- A flat stone. Opens the Horse Cemetery gate.
  \hfill \textit{The gate is a skull. The key fits the eye socket. The skull will scream. The scream is a name — the name of the horse you must bury.} `[MAGIC]`
\item \textbf{Iron Corral Nail} --- One sword blade from the fence. Use it to seal an oath.
  \hfill \textit{The blade is rusty. The rust is from broken promises. It will rust your oath too. Speak true, or bleed.} `[OATH]`
\item \textbf{Sky-Road Stone (Chip)} --- A piece of the blue cairn. Carry it to avoid getting lost on the steppe.
  \hfill \textit{The stone is warm. It points to the north. The north is where the grass is greenest. The stone hums at dawn.} `[TRAVEL]`
\item \textbf{Khagan's Whistle (Horse Phalange)} --- A brass whistle. Blow it to summon a remount from any station.
  \hfill \textit{The whistle is a horse's toe bone. The horse was the Khagan's favourite. It will come — but only once.} `[MAGIC]`
\item[J] \textbf{Ghost Herd's Bridle (Leather)} --- A strap of old leather. Use it to command one spirit horse.
  \hfill \textit{The bridle is from a horse that drowned in the salt pan. The spirit is still wet. It will carry you across water.} `[DEATH]`
\item[Q] \textbf{Wind-Singer's Staff (Tent Pole)} --- A wooden rod. Plant it to stop the wind for one hour.
  \hfill \textit{The staff was a tent pole. The tent was the singer's home. The home is gone. The wind will listen — for one hour.} `[WEATHER]`
\item[K] \textbf{Khagan's Saddle (Felt Pad)} --- Ride in it to gain +1 Position on all mounted rolls.
  \hfill \textit{The saddle has no stirrups. You must learn to ride without them. The Khagan did. His ghost rides behind you.} `[COMBAT]`
\item[A] \textbf{Steppe's Edge Stone} --- A piece of the boundary where grass meets desert. Touch it to make the grass grow.
  \hfill \textit{The stone is a tooth. The tooth is from the Steppe itself. The Steppe is hungry. The grass will grow — but it will be thirsty.} `[LEGEND]`
\end{enumerate}

\section*{Quick use notes}
\label{sec:kuvani-quick-use}
\begin{itemize}
\item Draw until you have all four suits: Spade = place, Heart = actor, Club = pressure, Diamond = leverage. Highest rank sets the main Clock (2–5 → 4, 6–10 → 6, J/Q/K → 8, A → 10).
\item Diamonds are codified outcomes (marks, compacts, flutes) that change position rather than call for a roll.
\item If any A appears, echo \textbf{sky-omens}—hawks that circle counter‑clockwise, grass that bends against the wind, a bell that rings without being struck, a shadow that falls where no sun shines.
\end{itemize}

\subsection*{Additional Features (Cultural Mechanics)}
\label{sec:kuvani-mechanics}

\begin{itemize}
\item \textbf{The Story Toll:} Kuvani remount stations do not accept coin. To take a fresh horse, you must tell a story that the remount master has not heard. The story is judged by the horse — if the horse stamps its foot, the story is false. If it whinnies, the story is true. If it remains still, the story is boring — you walk. Kuva's hawks watch; a hawk's cry means the story is worthy of the sky.
\item \textbf{Wind-Songs and Sky-Roads:} The Kuvani navigate by the stars and the wind. A Wind-Singer can hum a melody that encodes the route to water, to grass, or to safety. The song changes with the season. Only the singer knows the current melody. Those who hear the song without permission are cursed — they will hear it in their dreams until they pay a horse to the singer.
\item \textbf{Grazing Compacts as Currency:} A grazing compact (a sheepskin document sealed with horsehair) can be traded, sold, or used as proof of passage. A compact from a powerful clan is worth more than a chest of silver. Forging a compact is punishable by exile — the forger is left on the steppe without a horse. Kuva will not watch over them.
\item \textbf{Kuva's Witness (Oath Mechanic):} Any oath sworn under open sky (no roof, no canopy) is witnessed by Kuva. Breaking such an oath imposes −1 die on all rolls until the oath is fulfilled or atoned for by a sacrifice: a horse, a memory, or a year of service to the wronged party. The GM may also tick a \textbf{Kuva's Attention [6]} clock; when full, a sky-omen (drought, storm, ghost herd) strikes the oath‑breaker.
\item \textbf{Local Patrons (Syncretic Names):}
  \begin{itemize}
    \item \textbf{Kuva} — The Eternal Blue Sky. *The witness of all oaths, the father of horses, the wind that carries prayers.*
    \item \textbf{The White Mare} — The Traveler. *The first horse, who showed the Kuvani the sky‑road. Her hoofprints are the stars.*
    \item \textbf{The Iron Hoof} — Khemesh. *The pressure beneath the steppe; the hunger that pulls horses down into the salt pan.*
    \item \textbf{The Shadow Khagan} — Malachai. *A cursed king who sold his clan's horses for immortality; he rides alone, and his horses have no shadows.*
  \end{itemize}
\end{itemize}

\subsection*{Lore Echoes (Cross-Expansion Hooks)}
\begin{itemize}
  \item \textbf{Remount Chains and the Caravan Trade (Way of Silk):} Every major caravan on the eastern route uses Kuvani remount stations. A disruption in the remount chain — a sickness, a feud, a stolen horse — can halt trade for an entire season. The Tulkani have a saying: ``A Kuvani horse is worth nine stories. A dead Kuvani horse is worth nine feuds.''
  \item \textbf{Grazing Compacts and the Ykrul:} The Ykrul are the Kuvani's rivals and sometimes allies. A grazing compact with a Kuvani clan may require a Ykrul hostage string as surety. Breaking the compact means the hostage is executed — and Kuva will turn his face from the oath‑breaker. The Ykrul do not fear Kuva, but they respect his silence.
  \item \textbf{Wind-Singers and the Sky (Kuva):} The Kuvani worship Kuva as the watcher of oaths. Their wind-songs are prayers to the sky. A Wind-Singer who loses her voice is considered cursed by Kuva; she must walk the steppe alone until she hears the sky again. Some never return; their bones become sky‑road markers.
  \item \textbf{The Ghost Herd and the Hollow (Eastern Realms):} The ghost herd at the salt pan is said to be the Hollow's own horses — ridden by the Oath-Keepers. Aelaerem leave offerings of salt to keep the herd from riding through their villages. The Kuvani avoid the salt pan entirely. They say the ghost herd is Kuva's punishment for a Khagan who broke an oath to the sky.
  \item \textbf{The Iron Corral and the Black Banners:} The sword blades of the Iron Corral are from a war that never ended — a condotta that was paid twice, first in coin, then in blood. The Black Banner captains know the story; they leave offerings of rust at the corral when they pass. A captain who refuses to pay finds his horses spooked and his supply lines cut.
\end{itemize}

\subsection*{Kuvani Superstitions (burn for +1 die once)}
\begin{itemize}
  \item \textbf{Bread to the Horses:} Crumble a crust into the corral. Ignore the first \emph{Horse Sickness} penalty. Kuva's hawks will watch but not strike.
  \item \textbf{Knot the Hostage String:} Tie a new knot in a borrowed string. Cancel \emph{Hostage String Cut} or take –1 die on your next negotiation with Ykrul (your choice). The knot must be tied while facing the rising sun.
  \item \textbf{Name to the Wind:} Whisper a dead horse's name into the grass. Gain +1 Effect when riding this scene, but mark \emph{Obligation} to the Wind-Singer. The wind will carry the name to Kuva. He may answer with a storm.
\end{itemize}

\subsection*{Thematic SB Spend Table}
\label{sec:kuvani-sb}

\paragraph{Minor Complications (1 SB)}
\begin{itemize}
\item \textbf{Exposure:} Your actions draw unwanted attention from \textbf{Remount Masters or Wind-Singers}.
\item \textbf{Noise:} Sounds of your actions alert nearby \textbf{herdsmen or felt village guards}.
\item \textbf{Trace:} Evidence of your passage marks your route for \textbf{rival clans or ghost herd scouts}.
\item \textbf{Delay:} A brief but meaningful setback costs you \textbf{time or favourable grazing access}.
\item \textbf{Supply Strain:} Mark +1 segment on a relevant \textbf{resource clock}.
\end{itemize}

\paragraph{Moderate Setbacks (2 SB)}
\begin{itemize}
\item \textbf{Alarm Raised:} \textbf{Local Remount Master or Wind-Singer} becomes aware and begins responding.
\item \textbf{Position Lost:} You lose advantageous ground/cover/stealth due to \textbf{grass fire or sudden fog}.
\item \textbf{Foe Appears:} A \textbf{rival clan rider or Ykrul envoy} arrives on scene.
\item \textbf{Gear Trouble:} A piece of equipment becomes \textbf{Compromised/Neglected}.
\item \textbf{Lock/Barrier:} A simple obstacle now requires a test to overcome.
\end{itemize}

\paragraph{Serious Trouble (3 SB)}
\begin{itemize}
\item \textbf{Reinforcements:} Additional \textbf{Kuvani riders, Ykrul warriors, or ghost herd} arrive.
\item \textbf{Key Gear Breaks:} A crucial tool/weapon becomes temporarily unusable.
\item \textbf{Major Twist:} The situation fundamentally changes—\textbf{water poisoned/hostage string cut/khagan's demand issued}.
\item \textbf{Rail Tick:} Advance a relevant campaign/front clock by 1 segment.
\item \textbf{Condition Applied:} Mark \textbf{Fatigue 1/Harm 1/Condition} appropriate to fiction.
\end{itemize}

\paragraph{Major Turns (4+ SB)}
\begin{itemize}
\item \textbf{Trap Springs:} A prepared danger activates with full effect.
\item \textbf{Authority Arrival:} \textbf{Khagan, Wind-Singer of the White Mane, or Ghost Herd's Rider} intervenes.
\item \textbf{Scene Shift:} The environment changes dramatically—\textbf{steppe catches fire/salt pan floods/sky-road marker falls}.
\item \textbf{Patron Omen:} Divine/arcane forces take notice—\textbf{omen appears/blessing lost/curse manifests}.
\item \textbf{Narrative Pivot:} The story takes an unexpected turn that reframes objectives.
\end{itemize}

\paragraph{Region-Specific SB Options}
\begin{itemize}
\item \textbf{Kuvani (Remount Law):} Horse bells ring out of sequence, saddle leather creaks a warning, a horse refuses to move.
\item \textbf{Kuvani (Wind-Song):} Grass whispers a name, the wind stops mid-sentence, a bone flute plays itself.
\item \textbf{Kuvani (Ghost Herd):} Hoofprints appear where there are no horses, salt tastes of blood, a skeletal horse watches from the horizon.
\end{itemize}

\subsection*{Pursuit Ladder (Near / Mid / Far)}
\begin{itemize}
\item \textbf{Advance/Withdraw (Wits + Ride):} On a hit, shift one band. On a strong hit, also impose \emph{Wind at Back}.
\item \textbf{Cut the Herd (Presence + Sway):} On a hit, freeze range for one exchange; if a remount chain mark is in your possession, +1 die.
\item \textbf{Weather the Curse (Resolve + Spirit):} On a hit in \emph{Horse Sickness}, you pick which mount survives; on a miss, the GM does.
\end{itemize}

\subsection*{Boss Seeds (Hoof \& Wind)}
\begin{itemize}
\item \textbf{The Bone-Singer (Nine Plagues)} — \emph{Phases}: Horse Sickness Spread $\rightarrow$ Orchestrated Famine $\rightarrow$ Herd Coup. \emph{Cracks}: find a true story the bone-singer has not heard, expose a false cure, survive a stampede. \emph{Kuva's Judgment}: if the party swears an oath to end the sickness under open sky, the bone-singer loses one phase.
\item \textbf{The Storm Khagan (Rider of the White Blizzard)} — \emph{Phases}: Grass Fires Called Mercy $\rightarrow$ Procession of Empty Saddles $\rightarrow$ Wind-Baptism of a Remount Master. \emph{Cracks}: relic wind-song, survivor testimony, the sky-road turning against them. \emph{Kuva's Favour}: if a party member offers their own horse to the sky (releasing it onto the steppe), the Storm Khagan must pause for one exchange.
\end{itemize}

\begin{tcolorbox}[colback=black!3,colframe=black!40!white,title={Patronage \& Power}]
Among the Kuvani, power flows through remount chains, wind-songs, and the truth of a story. A clan's authority rests on the speed of its horses and the skill of its singers. The Khagan is a figurehead; the real power lies with the Remount Masters who control the stations, and the Wind-Singers who read the grass — and Kuva's will.

\textbf{For the GM:}  
Power in Kuvani society is measured in horse teeth, sheepskin compacts, and the memory of a good story. Patronage often takes the form of remount marks, grazing compacts, or bone flutes — each tied to a specific station, clan, or singer. To emphasize the itinerant and sky‑worshipping culture:
\begin{itemize}
\item Require players to \textbf{tell a story} to access remounts, shelter, or grazing rights. The GM can judge the story's quality (or have the player roll Presence + Sway vs. DV based on the remount master's mood). A failed story means the player must offer a truth — a secret, a confession, or a memory — to Kuva directly (speak it aloud under open sky).
\item Use \textbf{Wind-Singers} as quest-givers and as intermediaries with Kuva. They hear songs on the wind that only the party can resolve. The song is always a warning or a debt call, and often a prophecy.
\item Let \textbf{grazing compacts} be a form of currency that can be stolen, forged, or voided. A party without a compact is trespassing — the clans can demand a story as payment, or seize a horse as collateral. A compact sealed with a hair from the Khagan's mare is unbreakable without divine intervention.
\item Emphasize that \textbf{the steppe is alive and Kuva watches}: the grass remembers, the wind listens, and a broken promise is written in the sky-road. The Hollow's attention ticks when the party lies to a Kuvani. Kuva's attention ticks when they break an oath sworn under open sky.
\item Use \textbf{sky-omens} as scene-setting tools: a hawk's cry, a sudden stillness, a cloud shaped like a horse. These can foreshadow complications or reward truthful behaviour.
\end{itemize}
In Kuvani society, the grass forgets nothing — and Kuva's sky remembers every word spoken in the open.
\end{tcolorbox}

% Index entries
\index{Kuvani|see{Steppe Hosts}}
\index{Theme generator}
\index{Steppe Hosts generator}
\index{Places!Kuvani}
\index{People!Kuvani}
\index{Complications!Kuvani}
\index{Rewards!Kuvani}
\index{Remount stations}
\index{Wind-Singers}
\index{Grazing compacts}
\index{Horse sickness}
\index{Ghost herd}
\index{Sky-road}
\index{Iron corral}
\index{Khagan}
\index{Ykrul}
\index{Story toll}
\index{Kuva}
\index{Sky deity}
\index{Local Patrons!Kuvani}